% META: {"id": "1NSI-WEB-EVAL-B-corrige", "chapitre": "1NSI-WEB-IHM", "type_objet": "corrige_evaluation", "evaluation_ref": "1NSI-WEB-EVAL-B", "capacites": ["P-IHM-01A", "P-IHM-01B", "P-IHM-02", "P-IHM-03A", "P-IHM-03B", "P-IHM-03C", "P-IHM-04A", "P-IHM-04B", "P-IHM-04C"], "mode_creation": "ex_nihilo", "sources_inspiration": [], "status": "verified"}
% BEGIN-VERIFY
% gestionnaires = {}
% def ajouter_gestionnaire(id_composant, fonction):
%     gestionnaires[id_composant] = fonction
% def declencher_clic(id_composant):
%     if id_composant in gestionnaires:
%         return gestionnaires[id_composant]()
%     return None
% favoris = {"articles": 3}
% def retirer_favori():
%     favoris["articles"] -= 1
%     return favoris["articles"]
% ajouter_gestionnaire("bouton_retirer", retirer_favori)
% resultats = [declencher_clic("bouton_retirer") for _ in range(2)]
% assert resultats == [2, 1]
% from urllib.parse import urlencode
% params = {"ville": "Sousse", "type": "appartement"}
% url = "https://immo.fr/recherche?" + urlencode(params)
% assert url == "https://immo.fr/recherche?ville=Sousse&type=appartement"
% gestionnaires = {}
% journal = []
% def ajouter_gestionnaire(id_composant, evenement, fonction):
%     gestionnaires[(id_composant, evenement)] = fonction
% def declencher(id_composant, evenement, *args):
%     cle = (id_composant, evenement)
%     if cle in gestionnaires:
%         return gestionnaires[cle](*args)
%     return None
% ajouter_gestionnaire("recherche", "saisie", lambda texte: journal.append(("filtre", texte)))
% ajouter_gestionnaire("tri", "changement", lambda choix: journal.append(("tri", choix)))
% ajouter_gestionnaire("valider", "clic", lambda: journal.append(("envoi", None)))
% declencher("recherche", "saisie", "alg")
% declencher("tri", "changement", "prix")
% declencher("valider", "clic")
% assert journal == [("filtre", "alg"), ("tri", "prix"), ("envoi", None)]
% assert declencher("recherche", "clic") is None
% journal.clear()
% for lettre in "abc":
%     declencher("recherche", "saisie", lettre)
% declencher("valider", "clic")
% assert [e[0] for e in journal] == ["filtre", "filtre", "filtre", "envoi"]
% session = {}
% def connexion(identifiant, mot_de_passe, chiffre):
%     assert chiffre, "le mot de passe ne doit pas circuler en clair"
%     session["jeton"] = "jeton-" + identifiant
%     return session["jeton"]
% connexion("eleve", "secret", chiffre=True)
% assert session == {"jeton": "jeton-eleve"}
% assert "secret" not in session
% try:
%     connexion("eleve", "secret", chiffre=False)
% except AssertionError as e:
%     assert "clair" in str(e)
% else:
%     raise AssertionError("une connexion non chiffree aurait du etre refusee")
% END-VERIFY
\begin{corrige}{1NSI-WEB-EVAL-B-EX1}
\textbf{1.} Non. Le HTML décrit uniquement la structure et le contenu des composants ; sans code
JavaScript associant une fonction à l'événement clic, rien ne se produit automatiquement.

\textbf{2.}
\begin{python}
favoris = {"articles": 3}

def retirer_favori():
    favoris["articles"] -= 1
    return favoris["articles"]

ajouter_gestionnaire("bouton_retirer", retirer_favori)
\end{python}

\textbf{3.} La séquence est \lstinline{[2, 1]}.
\end{corrige}

\begin{corrige}{1NSI-WEB-EVAL-B-EX2}
\textbf{1.}
\begin{python}
from urllib.parse import urlencode

params = {"ville": "Sousse", "type": "appartement"}
url = "https://immo.fr/recherche?" + urlencode(params)
print(url)
# https://immo.fr/recherche?ville=Sousse&type=appartement
\end{python}

\textbf{2.} Oui, GET est adapté : les critères de recherche (ville, type de bien) ne sont
pas confidentiels, et l'URL résultante peut être partagée ou mise en favori.

\textbf{3.} Le formulaire comporte trois champs saisis par le visiteur : le nom, l'email
et le message. Il faut privilégier \textbf{POST} : ces données à caractère personnel ne
doivent pas apparaître dans l'URL, où elles seraient visibles, mises en favori et
enregistrées dans l'historique et les journaux du serveur. Avec POST, le serveur reçoit
une requête dont le \textbf{corps} contient les trois couples champ/valeur ; l'URL, elle,
ne porte que le chemin du script de traitement.
\end{corrige}

\begin{corrige}{1NSI-WEB-EVAL-B-EX3}
\textbf{1.} Trois associations possibles :
\begin{itemize}
  \item champ de texte \lstinline{"recherche"} : evenement de \textbf{saisie}, declenche a
  chaque frappe -- on filtre la liste affichee au fur et a mesure ;
  \item menu deroulant \lstinline{"tri"} : evenement de \textbf{changement de selection} --
  on reordonne la liste selon le critere choisi ;
  \item bouton \lstinline{"valider"} : evenement de \textbf{clic} -- on envoie la requete au
  serveur.
\end{itemize}
Sans fonction associee, le composant reste inerte : le HTML decrit l'apparence, il ne decrit
aucun traitement.

\textbf{2.} L'evenement de saisie se declenche \textbf{plusieurs fois}, a chaque caractere,
et sert a un retour immediat sans quitter la page ; l'evenement de validation se declenche
\textbf{une seule fois} et engage une action couteuse ou definitive, comme l'envoi au
serveur. Confondre les deux ferait envoyer une requete a chaque lettre tapee.
\end{corrige}

\begin{corrige}{1NSI-WEB-EVAL-B-EX4}
\textbf{1.} Dans l'ordre :
\begin{itemize}
  \item \textbf{Authentification} : le \textbf{client} affiche le formulaire et verifie
  eventuellement que les champs sont remplis ; le \textbf{serveur} verifie l'identifiant et
  le mot de passe, car lui seul detient les donnees de reference.
  \item \textbf{Affichage des places} : le \textbf{serveur} consulte la base et envoie
  l'etat des places ; le \textbf{client} met en forme et gere les interactions locales
  (survol, selection).
  \item \textbf{Enregistrement} : le \textbf{client} envoie la place choisie ; le
  \textbf{serveur} enregistre la reservation -- lui seul peut garantir qu'une place n'est pas
  attribuee deux fois.
\end{itemize}

\textbf{2.} Le navigateur memorise un \textbf{identifiant de session} (cookie ou jeton),
ainsi que l'etat d'affichage courant. C'est cet identifiant qui est \textbf{retransmis a
chaque requete} : le mot de passe, lui, ne circule qu'une fois, a la connexion, et n'est pas
conserve cote client.

\textbf{3.} Le chiffrement est indispensable \textbf{des la premiere etape}, celle qui
transporte le mot de passe, et il doit ensuite etre maintenu car le jeton de session
circule a chaque requete : intercepte, il vaudrait le mot de passe. Le chiffrement protege
la \textbf{confidentialite} et l'\textbf{integrite} de ce qui transite entre le navigateur
et le serveur ; il ne protege ni le poste de l'utilisateur, ni le serveur lui-meme.
\end{corrige}
